\section{Methodological Framework}
\label{sec:methodology}

Our framework processes heterogeneous 3D medical volumes through four sequential stages: spatial standardization, domain leakage prevention, model architecture exploration, and post-hoc probability calibration.

\subsection{Spatial \& Volumetric Standardization}
\label{sec:spatial}

Raw medical scans exhibit multi-axial anisotropy and varying field-of-view dimensions across acquisition centers. We apply a 3-step spatial standardization pipeline:
\begin{enumerate}
    \item \textbf{RAS Reorientation:} NIfTI affine matrices reorient all 3D volumes to standard Right-Anterior-Superior (RAS) anatomical coordinates.
    \item \textbf{Isotropic Resampling:} Trilinear interpolation resamples non-uniform voxels ($1.47\text{--}3.90~\text{mm}^3$) to a uniform $2.0~\text{mm}^3$ isotropic resolution.
    \item \textbf{Center-of-Mass Grid Alignment:} The spatial Center of Mass ($\text{COM}$) of thresholded non-air voxels ($>5\text{th}$ percentile) is translated to the grid center ($C = 39.5$) of an $80 \times 80 \times 80$ matrix.
\end{enumerate}
Per-volume Z-score normalization ($\tilde{v} = (v - \mu_v)/(\sigma_v + 10^{-8})$) is computed independently per scan to prevent dataset-level summary statistics from leaking across folds.

\subsection{Preventing Scanner-Site Data Leakage}
\label{sec:leakage}

To prevent models from exploiting scanner-specific noise, we enforce \textbf{Site-Aware Stratified Group K-Fold Cross-Validation} ($K=5$). Let $S_i \in \{1, \dots, 15\}$ denote the scanner site ID for sample $i$, and $y_i \in \{0, 1\}$ denote the label. Folds are constructed such that:
\begin{equation}
S_{\text{train}}^{(k)} \cap S_{\text{val}}^{(k)} = \emptyset, \quad \forall k \in \{1, \dots, 5\}
\end{equation}
This ensures that validation scans originate strictly from acquisition centers unseen during training, providing an un-biased estimate of clinical generalizability.

\subsection{3D Deep ResNet Architectures}
\label{sec:arch}

We evaluate two 3D Deep ResNet architectures designed for $80^3$ volumetric inputs:
\begin{itemize}
    \item \textbf{Net3dR (845K params):} Initial $3\times3\times3$ Conv (16 channels) $\to$ MaxPool $\to$ 3 Residual Blocks $\to$ 3 Downsampling Blocks (channel widths: 16--32--64--128) $\to$ AdaptiveAvgPool $\to$ FC Classifier.
    \item \textbf{Net3dBig (1.3M params):} Wider channel progression (20--40--80--160) with 3 Residual Blocks and Dropout ($p=0.4$).
\end{itemize}

Models are trained using Binary Cross-Entropy with Mixup augmentation ($\alpha=0.2$) and Label Smoothing ($\epsilon=0.05$):
\begin{equation}
y_{\text{smooth}} = y(1 - \epsilon) + 0.5\epsilon
\end{equation}
Optimization uses AdamW ($\text{lr}=10^{-3}$, weight decay $10^{-4}$) with Cosine Annealing learning rate schedules and gradient accumulation (effective batch size 24).

\subsection{Multi-View Test-Time Augmentation (TTA)}
\label{sec:tta}

During inference, each volume undergoes 15-view TTA: 1 base view, 6 spatial translation views ($\pm 1$ voxel along $X, Y, Z$), and 8 axial/sagittal rotation views ($\pm 2^\circ, \pm 4^\circ$). Predictions are averaged across all models and views.

\subsection{Post-Hoc Probability Calibration}
\label{sec:calibration}

Raw predictions $p \in (0, 1)$ are calibrated using Platt Scaling~\cite{platt1999probabilistic} fit on out-of-fold validation logits $z = \text{logit}(p)$:
\begin{equation}
p_{\text{cal}} = \frac{1}{1 + \exp(-(a \cdot z + b))}
\end{equation}
Parameters $(a=1.45, b=0.20)$ are optimized to minimize out-of-sample LogLoss and Expected Calibration Error (ECE).